\subsection{Matrix-Induced Maps and Rank}
\begin{definition} \hlt{Linear Map Induced by a Matrix}\\
For $A\in\operatorname{Mat}(m,n,\mathbb{F})$, the \hlt{linear map induced by $A$} is the map $T_A:\mathbb{F}^n\to\mathbb{F}^m$ defined by $T_A(\mathbf{x})=A\mathbf{x}$.
\end{definition}

\begin{proposition} \label{prop:matrix-induced-map-linear} \hlt{Matrix-Induced Maps Are Linear}\\
If $A\in\operatorname{Mat}(m,n,\mathbb{F})$, then $T_A:\mathbb{F}^n\to\mathbb{F}^m$ is linear.
\begin{proof}
For $\mathbf{x},\mathbf{y}\in\mathbb{F}^n$ and $a\in\mathbb{F}$, $T_A(a\mathbf{x}+\mathbf{y})=A(a\mathbf{x}+\mathbf{y})=aA\mathbf{x}+A\mathbf{y}=aT_A(\mathbf{x})+T_A(\mathbf{y})$.\\
\hlt{Structural Commentary}: Matrix multiplication preserves addition and scalar multiplication.\\
\hlt{Key Insight}: Every matrix defines a canonical linear map between coordinate spaces.
\end{proof}
\end{proposition}

\begin{corollary} \label{cor:matrix-rank-nullity} \hlt{Matrix Rank-Nullity}\\
Let $A\in\operatorname{Mat}(m,n,\mathbb{F})$. Then $\dim(\ker(A))+\dim(\operatorname{col}(A))=n$ and $\operatorname{nullity}(A)+\operatorname{rank}(A)=n$.\\
If $m=n$, then $\ker(A)=\{\mathbf{0}\}$ if and only if $\operatorname{col}(A)=\mathbb{F}^n$.
\begin{proof}
Apply Theorem \ref{thm:rank-nullity} to the induced linear map $T_A:\mathbb{F}^n\to\mathbb{F}^m$. By definition, $\ker(T_A)=\ker(A)$ and $\operatorname{range}(T_A)=\operatorname{col}(A)$, giving the displayed identities. When $m=n$, the remaining equivalence follows as:
\begin{enumerate}[label=\roman*.]
\setlength{\itemsep}{0pt}
\item[$(\Rightarrow)$] If $\ker(A)=\{\mathbf{0}\}$, then $\dim(\operatorname{col}(A))=n$. Since $\operatorname{col}(A)\subseteq\mathbb{F}^n$, Proposition \ref{prop:dimension-bounds-and-basis-criteria} gives $\operatorname{col}(A)=\mathbb{F}^n$.
\item[$(\Leftarrow)$] If $\operatorname{col}(A)=\mathbb{F}^n$, then $\dim(\operatorname{col}(A))=n$, so the identity gives $\dim(\ker(A))=0$. Hence $\ker(A)=\{\mathbf{0}\}$.
\end{enumerate}
\hlt{Structural Commentary}: Matrix rank-nullity is ordinary rank-nullity written in coordinates.\\
\hlt{Key Insight}: Kernel loss plus column-space dimension equals the number of columns.
\end{proof}
\end{corollary}

\begin{theorem} \label{thm:row-rank-equals-column-rank} \hlt{Row Rank Equals Column Rank}\\
For $A\in\operatorname{Mat}(m,n,\mathbb{F})$, $\dim(\operatorname{row}(A))=\dim(\operatorname{col}(A))$. Equivalently, $\operatorname{rank}(A^t)=\operatorname{rank}(A)$.
\begin{proof}
Let $r=\operatorname{rank}(A)$. If $r=0$, then $\operatorname{col}(A)=\{\mathbf{0}\}$, so every column of $A$ is zero and $A=0$; the result follows. Suppose $r>0$. Choose columns $\mathbf{a}_{j_1},\ldots,\mathbf{a}_{j_r}$ forming a basis of $\operatorname{col}(A)$, and set $B=[\mathbf{a}_{j_1}\ \cdots\ \mathbf{a}_{j_r}]\in\operatorname{Mat}(m,r,\mathbb{F})$.\\
For each column $\mathbf{a}_k$ of $A$, there is a unique $\mathbf{x}_k\in\mathbb{F}^r$ such that $\mathbf{a}_k=B\mathbf{x}_k$. With $X=[\mathbf{x}_1\ \cdots\ \mathbf{x}_n]\in\operatorname{Mat}(r,n,\mathbb{F})$, we have $A=BX$. Therefore $A^t=X^tB^t$, so every column of $A^t$ is a linear combination of the $r$ columns of $X^t$. Hence $\operatorname{rank}(A^t)\leq r=\operatorname{rank}(A)$. Applying the same argument to $A^t$ gives $\operatorname{rank}(A)\leq\operatorname{rank}(A^t)$.\\
\hlt{Structural Commentary}: Column rank controls row rank by factoring $A$ through a basis of its column space.\\
\hlt{Key Insight}: Rows and columns measure the same intrinsic rank, viewed from transposed coordinates.
\end{proof}
\end{theorem}

\begin{perspective} \label{per:row-column-rank-different-hosts} \hlt{Row and Column Spaces Have Different Hosts}\\
For $A\in\operatorname{Mat}(m,n,\mathbb{F})$, $\operatorname{row}(A)\subseteq\mathbb{F}^n$ and $\operatorname{col}(A)\subseteq\mathbb{F}^m$.\\
Hence, when $m\neq n$, Theorem \ref{thm:row-rank-equals-column-rank} asserts only $\dim\operatorname{row}(A)=\dim\operatorname{col}(A)$, not equality of subspaces.
\end{perspective}

\begin{theorem} \label{thm:matrix-four-fundamental-subspaces-dimensions} \hlt{Four Fundamental Matrix Subspaces}\\
Let $A\in\operatorname{Mat}(m,n,\mathbb{F})$ and let $r=\operatorname{rank}(A)$. Then $\operatorname{row}(A),\ker(A) \subseteq \mathbb{F}^n$; $\operatorname{col}(A),\ker(A^t) \subseteq \mathbb{F}^m$,\\
and $\dim(\operatorname{col}(A))=\dim(\operatorname{row}(A))=r$; $\dim(\ker(A))=n-r$; $\dim(\ker(A^t))=m-r$.
\begin{proof}
The inclusions follow from the ambient dimensions of the rows and columns, and from the domains of $A$ and $A^t$. By definition, $\dim(\operatorname{col}(A))=r$; by Theorem \ref{thm:row-rank-equals-column-rank}, $\dim(\operatorname{row}(A))=r$.\\
Rank-nullity gives $\dim(\ker(A))=n-\operatorname{rank}(A)=n-r$. Since $\operatorname{rank}(A^t)=\dim(\operatorname{col}(A^t))=\dim(\operatorname{row}(A))=r$, rank-nullity for $A^t:\mathbb{F}^m\to\mathbb{F}^n$ gives $\dim(\ker(A^t))=m-\operatorname{rank}(A^t)=m-r$.\\
\hlt{Structural Commentary}: One rank determines all four fundamental dimensions.\\
\hlt{Key Insight}: Rank counts what survives; nullity counts what is killed.
\end{proof}
\end{theorem}

\begin{figure}[H]
  \centering
  \scalebox{0.9}{\begin{tikzpicture}[
  scale=0.90,
  transform shape,
  >=Stealth,
  every node/.style={font=\small},
  host/.style={draw=gray!50, fill=gray!3, rounded corners=3pt, line width=0.55pt,
    minimum width=4.05cm, minimum height=2.05cm},
  entry/.style={draw=gray!60, rounded corners=2pt, line width=0.5pt,
    minimum width=1.72cm, minimum height=0.78cm, align=center},
  active/.style={entry, fill=teal!8},
  null/.style={entry, fill=blue!7},
  map/.style={->, gray!65!black, line width=0.75pt}
]
  \node[host] (Fn) at (0,0) {};
  \node[anchor=north west] at (-1.82,0.88) {$\mathbb{F}^n$};
  \node[host] (Fm) at (5.6,0) {};
  \node[anchor=north west] at (3.78,0.88) {$\mathbb{F}^m$};

  \node[null] at (-0.98,-0.05) {$\ker(A)$\\[-1pt]{\footnotesize $n-r$}};
  \node[active] at (0.98,-0.05) {$\operatorname{row}(A)$\\[-1pt]{\footnotesize $r$}};
  \node[active] at (4.62,-0.05) {$\operatorname{col}(A)$\\[-1pt]{\footnotesize $r$}};
  \node[null] at (6.58,-0.05) {$\ker(A^t)$\\[-1pt]{\footnotesize $m-r$}};

  \draw[map] (2.1,0.4) -- (3.5,0.4) node[midway, above] {$A$};
  \draw[map] (3.5,-0.4) -- (2.1,-0.4) node[midway, below] {$A^t$};
\end{tikzpicture}
}
  \caption{Four fundamental subspaces by host coordinate space.}
  \label{fig:four-fundamental-subspaces-hosts}
\end{figure}

\begin{corollary} \label{cor:homogeneous-system-more-variables-than-equations} \hlt{Homogeneous Systems with More Variables Than Equations}\\
If $A\in\operatorname{Mat}(m,n,\mathbb{F})$ with $n>m$, then the homogeneous system $A\mathbf{x}=\mathbf{0}$ has a nonzero solution.
\begin{proof}
The induced map $T_A:\mathbb{F}^n\to\mathbb{F}^m$ has domain dimension larger than codomain dimension, so Corollary \ref{cor:map-to-smaller-space-not-injective} shows that $T_A$ is not injective.\\
Hence $\ker(A)=\ker(T_A)$ contains a nonzero vector $\mathbf{x}$, and this vector satisfies $A\mathbf{x}=\mathbf{0}$.\\
\hlt{Structural Commentary}: Extra variables force a nontrivial kernel.\\
\hlt{Key Insight}: A nontrivial kernel is exactly a nonzero homogeneous solution.
\end{proof}
\end{corollary}

\begin{corollary} \label{cor:inhomogeneous-system-more-equations-than-variables} \hlt{Inhomogeneous Systems with More Equations Than Variables}\\
If $A\in\operatorname{Mat}(m,n,\mathbb{F})$ with $m>n$, then there exists $\mathbf{b}\in\mathbb{F}^m$ such that $A\mathbf{x}=\mathbf{b}$ has no solution.
\begin{proof}
The induced map $T_A:\mathbb{F}^n\to\mathbb{F}^m$ has domain dimension smaller than codomain dimension, so Corollary \ref{cor:map-to-larger-space-not-surjective} shows that $T_A$ is not surjective. Choose $\mathbf{b}\in\mathbb{F}^m\setminus\operatorname{range}(T_A)$. Then no $\mathbf{x}\in\mathbb{F}^n$ satisfies $A\mathbf{x}=\mathbf{b}$.\\
\hlt{Structural Commentary}: Too few variables cannot reach every possible right-hand side.\\
\hlt{Key Insight}: A missing range vector becomes an inconsistent inhomogeneous system.
\end{proof}
\end{corollary}

\begin{proposition} \label{prop:matrix-representation-action} \hlt{Matrix Representation Acts on Coordinates}\\
Let $\beta=(\mathbf{v}_1,\ldots,\mathbf{v}_n)$ be a basis of $\mathcal{V}$, let $\gamma$ be a basis of $\mathcal{W}$, and let $T\in\LL(\mathcal{V},\mathcal{W})$.\\
Then for every $\mathbf{v} \in \mathcal{V}$, $[T(\mathbf{v})]_{\gamma} = [T]_{\gamma \leftarrow \beta}[\mathbf{v}]_{\beta}$.
\begin{proof}
Write $\mathbf{v} = a_1\mathbf{v}_1 + \cdots + a_n\mathbf{v}_n$. Then linearity gives $T(\mathbf{v}) = a_1T(\mathbf{v}_1) + \cdots + a_nT(\mathbf{v}_n)$ and the right-hand side is exactly the linear combination of the columns of $[T]_{\gamma \leftarrow \beta}$ with coefficients $a_1,\ldots,a_n$.\\
\hlt{Structural Commentary}: Matrix representation turns applying $T$ into multiplying coordinate columns.\\
\hlt{Key Insight}: The columns of $[T]_{\gamma \leftarrow \beta}$ are the coordinate images of the basis vectors.
\end{proof}
\end{proposition}
